
\documentclass{article}
%%%%%%%%%%%%%%%%%%%%%%%%%%%%%%%%%%%%%%%%%%%%%%%%%%%%%%%%%%%%%%%%%%%%%%%%%%%%%%%%%%%%%%%%%%%%%%%%%%%%%%%%%%%%%%%%%%%%%%%%%%%%%%%%%%%%%%%%%%%%%%%%%%%%%%%%%%%%%%%%%%%%%%%%%%%%%%%%%%%%%%%%%%%%%%%%%%%%%%%%%%%%%%%%%%%%%%%%%%%%%%%%%%%%%%%%%%%%%%%%%%%%%%%%%%%%
\usepackage[left=1in,right=1in,top=1in,bottom=1in]{geometry}
\usepackage{graphicx}
\usepackage{float}
\usepackage{hyperref}

%TCIDATA{OutputFilter=LATEX.DLL}
%TCIDATA{Version=5.50.0.2960}
%TCIDATA{<META NAME="SaveForMode" CONTENT="1">}
%TCIDATA{BibliographyScheme=Manual}
%TCIDATA{Created=Tuesday, April 11, 2023 10:06:45}
%TCIDATA{LastRevised=Friday, January 24, 2025 17:25:22}
%TCIDATA{<META NAME="GraphicsSave" CONTENT="32">}
%TCIDATA{<META NAME="DocumentShell" CONTENT="Standard LaTeX\Blank - Standard LaTeX Article">}
%TCIDATA{CSTFile=40 LaTeX article.cst}

\newtheorem{theorem}{Theorem}
\newtheorem{acknowledgement}[theorem]{Acknowledgement}
\newtheorem{algorithm}[theorem]{Algorithm}
\newtheorem{axiom}[theorem]{Axiom}
\newtheorem{case}[theorem]{Case}
\newtheorem{claim}[theorem]{Claim}
\newtheorem{conclusion}[theorem]{Conclusion}
\newtheorem{condition}[theorem]{Condition}
\newtheorem{conjecture}[theorem]{Conjecture}
\newtheorem{corollary}[theorem]{Corollary}
\newtheorem{criterion}[theorem]{Criterion}
\newtheorem{definition}[theorem]{Definition}
\newtheorem{example}[theorem]{Example}
\newtheorem{exercise}[theorem]{Exercise}
\newtheorem{lemma}[theorem]{Lemma}
\newtheorem{notation}[theorem]{Notation}
\newtheorem{problem}[theorem]{Problem}
\newtheorem{proposition}[theorem]{Proposition}
\newtheorem{remark}[theorem]{Remark}
\newtheorem{solution}[theorem]{Solution}
\newtheorem{summary}[theorem]{Summary}
\newenvironment{proof}[1][Proof]{\noindent\textbf{#1.} }{\ \rule{0.5em}{0.5em}}
\newcommand{\pwd}{D:/DeepLearning/images_from_K/lesion_normal_comparison/report/}
\input{tcilatex}
\begin{document}

\title{Spatial Analysis}
\author{}
\maketitle

\section{Introduction}

We compare the ulcerative and intact regions of a piece of tissue as
follows. First, we run a custom pipeline (https://github.com/youyifong/tsp)
on an image to produce a catalog of cells in the tissue for each marker of
interest. The pipeline starts with image registration, which transforms
different marker images into one coordinate system so that we can compare
the distribution of different markers from the same piece of tissue. After
image registration, we perform cell segmenation using a Cellpose model
fine-tuned with a set of expert-annotated training images of immune cells in
skin tissues. For each cell thus identified, we compute the distance from
its center to the skin boundary and the DEJ\ line, whether it is located in
the dermis or epidermis and whether it is in the ulcerative or intact
region, and the intensity of staining. Given the cell catalogs, we run
statistical analyses in R (https://cran.r-project.org/) to compare the
distributions of immune cells between the ulcerative and intact regions
tabularly and graphically. Graphically, we represent the distribution of a
type of immune cells using smoothed scaled histogram
(https://github.com/youyifong/kyotil), in which the y-axis is the density of
immune cells, computed as the number of cells per 100x100 um2 tissue, and
the x-axis is the depth of tissue in100um intervals, either with reference
to the skin boundary or the DEJ\ line.

\bigskip

\section{Methods}

\bigskip Issues

\begin{itemize}
\item Need to review length

\begin{itemize}
\item For LS, epi is a lot wider than dermis. How do we draw the length?

\item DD intact, length correct? CD3 peak seems too high in Fig 1
\end{itemize}

\item For JM, region 2 is excluded from analysis.

\item Need to review segmentation. flow .5 seems to results in many more
cells than flow .4, e.g. for CD11c in DD
\end{itemize}

\bigskip

For each sample, we will compare the distribution of each markers in the
normal and lesion regions.

\begin{itemize}
\item Markers for DD: CD3, CD4, CD8a, CD11c, CD15, CD45, CD68, CD69, CD123,
CD279

\item Markers for CF: CD1a, CD3, CD4, CD8a, CD11c, CD15, CD45, CD69, CD123,
CD279 (CD68 did not pass QC)
\end{itemize}

For CD123, segmentation is done at flow=0.5 to get more masks in the bright
area and a threshold of 5000 is set for total intensity to remove masks in
the debris field in the lesion area.

\bigskip

To compute cell density, we use the size of biopsy at each depth.

We also use a rough, single number meaurement of the size of the skin tissue
lesion area in some tables, which is determined by morphology using wheat
germ agglutinin (WGA) flourescent stain. ImageJ is used to demarcate the
lesion area, then a straight line ROI is drawn to span from one side of the
lesion to the other side of the lesion, with the endpoints of the line
positioned at the outermost skin boundary. \ \ \ \ 

\bigskip

Some considerations on region boundary and distance metric: \ 

1. Should we draw the boundary between normal and lesion based on WGA only
or all markers?

2. Shall each ROI be convex?

For (1), the two choices lead are answering different questions. Drawing
boundary based on WGA only answers a stronger question than drawing the
boundary based on all markers.

For (2), a convex shape is a shape that a line connecting any two points in
the shape completely lies within the shape. Practically speaking, in this
context it is equivalent to drawing a straight line to separate normal from
lesion. Pros for straight boundary: the span between any cell and the
closest point on the skin lies completely within the region. So the distance
measure is well defined. In contrast, if a line from a cell to its closest
point on the skin surface cuts through both normal and lesion regions, the
interpretation of distance to skin surface is muddled. Pros for curved
boundary: Regions (which is a type of neighborhoods) do not necessarily have
geometric shapes.

These two questions are also intertwined. If we use WGA only to draw
boundaries, we can only use straight line boundaries because the difference
between normal and lesion is only visible in epidermis in WGA staining. A
straight line can be drawn in epidermis and extend to dermis, but not a
curve. \ \ \ \ \ \ \ \ \ \ \ \ \ \ \ 

%TCIMACRO{\TeXButton{clearpage}{\clearpage}}%
%BeginExpansion
\clearpage%
%EndExpansion

\section{Results}

%TCIMACRO{%
%\TeXButton{tab DD_Les_counts}{\begin{center}
%\begin{table}[h] \centering
%
%\input{\pwd tables/input/DD_Les_counts.tex}
%
%\caption{DD}
%\vspace{-5pt} 
%\label{tab DD counts}
%
%\end{table}
%\end{center}}}%
%BeginExpansion
\begin{center}
\begin{table}[h] \centering

\input{\pwd tables/input/DD_Les_counts.tex}

\caption{DD}
\vspace{-5pt} 
\label{tab DD counts}

\end{table}
\end{center}%
%EndExpansion

%TCIMACRO{%
%\TeXButton{tab CF_Les_counts.tex}{\begin{center}
%\begin{table}[h] \centering
%
%\input{\pwd tables/input/CF_Les_counts.tex}
%
%\caption{CF}
%\vspace{-5pt} 
%\label{tab DD counts}
%
%\end{table}
%\end{center}}}%
%BeginExpansion
\begin{center}
\begin{table}[h] \centering

\input{\pwd tables/input/CF_Les_counts.tex}

\caption{CF}
\vspace{-5pt} 
\label{tab DD counts}

\end{table}
\end{center}%
%EndExpansion

%TCIMACRO{%
%\TeXButton{tab ulcerative}{\begin{center}
%\begin{table}[h] \centering
%
%\input{\pwd tables/input/ulcerative_comparison.tex}
%
%\caption{Ulcerative.
%The cell counts are per mm2. The widths and seperation are measured in um.}
%\vspace{-5pt} 
%\label{tab ulcerative}
%
%\end{table}
%\end{center}}}%
%BeginExpansion
\begin{center}
\begin{table}[h] \centering

\input{\pwd tables/input/ulcerative_comparison.tex}

\caption{Ulcerative.
The cell counts are per mm2. The widths and seperation are measured in um.}
\vspace{-5pt} 
\label{tab ulcerative}

\end{table}
\end{center}%
%EndExpansion

%TCIMACRO{%
%\TeXButton{tab intact}{\begin{center}
%\begin{table}[h] \centering
%
%\input{\pwd tables/input/intact_comparison.tex}
%
%\caption{Intact.
%The cell counts are per mm2. The widths and seperation are measured in um.}
%\vspace{-5pt} 
%\label{tab intact}
%
%\end{table}
%\end{center}}}%
%BeginExpansion
\begin{center}
\begin{table}[h] \centering

\input{\pwd tables/input/intact_comparison.tex}

\caption{Intact.
The cell counts are per mm2. The widths and seperation are measured in um.}
\vspace{-5pt} 
\label{tab intact}

\end{table}
\end{center}%
%EndExpansion

%TCIMACRO{\TeXButton{clearpage}{\clearpage}}%
%BeginExpansion
\clearpage%
%EndExpansion

%TCIMACRO{%
%\TeXButton{fig:compare_DD_CF_JM}{\begin{center}
%\begin{figure}[H] \centering
%
%\includegraphics[width=\textwidth]{\pwd figures/compare_DD_CF_JM_dist2skin.png}\\
%\vspace{30pt}
%\includegraphics[width=\textwidth]{\pwd figures/compare_DD_CF_JM_dist2DEJ.png}
%
%\caption{Number of cells per 100x100 um2 tissue as a function of distance to DEJ or skin.
%The curves are smoothed over histograms and each bin in the histogram represents a 100 um interval. 
%Adding up the counts from all the bins gives the number of cells in tissue of 100 um in width.
%For example, CD123 epi count 1038/2862*100=36.
%}
%
%\vspace{-5pt}
%\label{fig:compare_DD_CF_JM}
%\end{figure}
%\end{center}}}%
%BeginExpansion
\begin{center}
\begin{figure}[H] \centering

\includegraphics[width=\textwidth]{\pwd figures/compare_DD_CF_JM_dist2skin.png}\\
\vspace{30pt}
\includegraphics[width=\textwidth]{\pwd figures/compare_DD_CF_JM_dist2DEJ.png}

\caption{Number of cells per 100x100 um2 tissue as a function of distance to DEJ or skin.
The curves are smoothed over histograms and each bin in the histogram represents a 100 um interval. 
Adding up the counts from all the bins gives the number of cells in tissue of 100 um in width.
For example, CD123 epi count 1038/2862*100=36.
}

\vspace{-5pt}
\label{fig:compare_DD_CF_JM}
\end{figure}
\end{center}%
%EndExpansion

%TCIMACRO{%
%\TeXButton{fig:compare_biopsies_intact}{\begin{center}
%
%\begin{figure}[H] \centering
%\includegraphics[width=\textwidth]{\pwd figures/compare_biopsies_intact_dist2skin.png}
%\caption{Number of cells per 100x100 um2 tissue as a function of distance to DEJ or skin.
%The curves are smoothed over histograms and each bin in the histogram represents a 100 um interval. 
%Adding up the counts from all the bins gives the number of cells in tissue of 100 um in width.
%}
%\end{figure}
%
%\vspace{30pt}
%
%\begin{figure}[H] \centering
%\includegraphics[width=\textwidth]{\pwd figures/compare_biopsies_intact_dist2DEJ.png}
%\caption{Number of cells per 100x100 um2 tissue as a function of distance to DEJ or skin.
%The curves are smoothed over histograms and each bin in the histogram represents a 100 um interval. 
%Adding up the counts from all the bins gives the number of cells in tissue of 100 um in width.
%}
%\end{figure}
%
%\end{center}}}%
%BeginExpansion
\begin{center}

\begin{figure}[H] \centering
\includegraphics[width=\textwidth]{\pwd figures/compare_biopsies_intact_dist2skin.png}
\caption{Number of cells per 100x100 um2 tissue as a function of distance to DEJ or skin.
The curves are smoothed over histograms and each bin in the histogram represents a 100 um interval. 
Adding up the counts from all the bins gives the number of cells in tissue of 100 um in width.
}
\end{figure}

\vspace{30pt}

\begin{figure}[H] \centering
\includegraphics[width=\textwidth]{\pwd figures/compare_biopsies_intact_dist2DEJ.png}
\caption{Number of cells per 100x100 um2 tissue as a function of distance to DEJ or skin.
The curves are smoothed over histograms and each bin in the histogram represents a 100 um interval. 
Adding up the counts from all the bins gives the number of cells in tissue of 100 um in width.
}
\end{figure}

\end{center}%
%EndExpansion

%TCIMACRO{%
%\TeXButton{fig compare_normal_lesion}{\begin{figure}[H] \centering
%\includegraphics[width=.95\textwidth]{\pwd figures/compare_normal_lesion_dist2skin.png}\\
%\vspace{0pt}
%\includegraphics[width=.95\textwidth]{\pwd figures/compare_normal_lesion_dist2DEJ.png}
%\caption{Number of cells per 100x100 um2 tissue as a function of distance to skin or DEJ.
%The curves are smoothed over histograms and each bin in the histogram represents a 100 um interval. 
%Adding up the counts from all the bins gives the number of cells in tissue of 100 um in width.
%}
%\end{figure}}}%
%BeginExpansion
\begin{figure}[H] \centering
\includegraphics[width=.95\textwidth]{\pwd figures/compare_normal_lesion_dist2skin.png}\\
\vspace{0pt}
\includegraphics[width=.95\textwidth]{\pwd figures/compare_normal_lesion_dist2DEJ.png}
\caption{Number of cells per 100x100 um2 tissue as a function of distance to skin or DEJ.
The curves are smoothed over histograms and each bin in the histogram represents a 100 um interval. 
Adding up the counts from all the bins gives the number of cells in tissue of 100 um in width.
}
\end{figure}%
%EndExpansion

Comparing the first two rows in the first column, we see that the ratio of
the area under the curve across both epi and dermis between red and
blue(ulcerative and intact) is about 1 for CD3/CF and 1/3 for CD3/DD. This
suggests that CD3 is depleted in lesion in DD, but not in CF.

Compare lesion with normal in CD15/CF, we see that the lesion peak is
shifted to the left compared to the normal peak. This pattern is more clear
in CF than in DD. We don't see this pattern for CD3.

%TCIMACRO{%
%\TeXButton{fig compare_markers.png}{\begin{center}
%\begin{figure}[H] \centering
%\includegraphics[height=\textheight]{\pwd figures/compare_markers.png}
%
%\caption{Number of cells per 100x100 um2 tissue as a function of distance to skin.
%The curves are smoothed over histograms and each bin in the histogram represents a 100 um interval. 
%Adding up the counts from all the bins gives the number of cells in tissue of 100 um in width.
%}
%
%\vspace{-5pt}
%\label{}
%\end{figure}
%\end{center}}}%
%BeginExpansion
\begin{center}
\begin{figure}[H] \centering
\includegraphics[height=\textheight]{\pwd figures/compare_markers.png}

\caption{Number of cells per 100x100 um2 tissue as a function of distance to skin.
The curves are smoothed over histograms and each bin in the histogram represents a 100 um interval. 
Adding up the counts from all the bins gives the number of cells in tissue of 100 um in width.
}

\vspace{-5pt}
\label{}
\end{figure}
\end{center}%
%EndExpansion

That the CD15 dist2skin peak in DD/lesion is near skin while the
corresponding dist2DEJ peak is near DEJ is due to the fact that the DEJ line
is very curvy.

%TCIMACRO{%
%\TeXButton{fig:CD15cyanCD3green}{\begin{center}
%\begin{figure}[H] \centering
%DD ($6540\times 5262$)\\
%\includegraphics[width=.9\textwidth]{\pwd DD_CD15cyan_CD3green.png}\\
%CF ($5296\times 4040$)\\
%\includegraphics[width=.9\textwidth]{\pwd CF_CD15cyan_CD3green.png}
%\caption{CD3: green, contrast limits 0-0.5; CD15: cyan. The intact/ulcerative boundary and the epidmeris/dermis boundary are shown in white. Outer boundaries are not tight.
%}
%\label{fig:CD15cyanCD3green}
%\end{figure}
%\end{center}}}%
%BeginExpansion
\begin{center}
\begin{figure}[H] \centering
DD ($6540\times 5262$)\\
\includegraphics[width=.9\textwidth]{\pwd DD_CD15cyan_CD3green.png}\\
CF ($5296\times 4040$)\\
\includegraphics[width=.9\textwidth]{\pwd CF_CD15cyan_CD3green.png}
\caption{CD3: green, contrast limits 0-0.5; CD15: cyan. The intact/ulcerative boundary and the epidmeris/dermis boundary are shown in white. Outer boundaries are not tight.
}
\label{fig:CD15cyanCD3green}
\end{figure}
\end{center}%
%EndExpansion

%TCIMACRO{%
%\TeXButton{fig:CD15cyanCD3green}{\begin{center}
%\begin{figure}[H] \centering
%JM ($6528\times 4008$)\\
%\includegraphics[width=.9\textwidth]{\pwd JM_CD15cyan_CD3green.png}
%\caption{CD3: green; CD15: cyan, contrast limits 0-0.5. 
%Region boudaries are shown in white. 
%}
%\label{fig:CD15cyanCD3green}
%\end{figure}
%\end{center}}}%
%BeginExpansion
\begin{center}
\begin{figure}[H] \centering
JM ($6528\times 4008$)\\
\includegraphics[width=.9\textwidth]{\pwd JM_CD15cyan_CD3green.png}
\caption{CD3: green; CD15: cyan, contrast limits 0-0.5. 
Region boudaries are shown in white. 
}
\label{fig:CD15cyanCD3green}
\end{figure}
\end{center}%
%EndExpansion

In Figure \ref{fig:CD15cyanCD3green}\ we see that DD lesion has less CD3 and
is further away from DEJ than CF. Correspondingly, in Figure \ref%
{fig:compare_DD_CF_JM}, we see that in the dermis region close to DEJ, the
count density is about 50 for CF/ulcerative, CF/intact, and CF/ulcerative;
but for DD/ulcerative, the count density is much lower.

Biospies with only intact regions seem to have lower number of immune cells.

\section{Discussion}

Neutraphil to lymphocyte ratio (NLR) in blood is a prognostic factor for
cancer. Lower is better.

\section{References}

Chan A, Bentzen S, Rout A, Miller K. Predicting if lung cancer will
relapse-the role of neutrophil/lymphocyte ratio. Hematol Oncol Stem Cell
Ther. 2023;16(2):158--61.

Costa T, Nogueiro J, Ribeiro D, Viegas P, Santos-Sousa H. Impact of serum
albumin concentration and neutrophil-lymphocyte ratio score on gastric
cancer prognosis. Langenbecks Arch Surg. 2023;408(1):57.

Chen W, Xin S, Xu B. Value research of NLR, PLR, and RDW in prognostic
assessment of patients with colorectal cancer. J Healthcare Eng.
2022;16(2022):7971415.

\end{document}
